\chapter{Section graded ring infrastructure: tensor powers and graded sections}
\label{chap:Picard_SectionGradedRing}

\providecommand{\PshMod}{\mathrm{PresheafOfModules}}
\providecommand{\ShMod}{\mathrm{SheafOfModules}}
\providecommand{\Modules}{\mathrm{Modules}}

The Hilbert-polynomial route of \cref{chap:Picard_QuotScheme} represents the
degree-\(m\) graded component of a fibre by the global sections
\(\Gamma(X_s, \mathcal{F}_s \otimes L_s^{\otimes m})\) of a twist of a coherent
sheaf by a tensor power of a line bundle, and it organises these components into
a graded ring \(R(X_s, L_s) = \bigoplus_{m \ge 0} \Gamma(X_s, L_s^{\otimes m})\)
(\cref{def:sectionGradedRing}) and a graded module
\(M(X_s, \mathcal{F}_s, L_s) = \bigoplus_{m \ge 0}
\Gamma(X_s, \mathcal{F}_s \otimes L_s^{\otimes m})\)
(\cref{def:sectionGradedModule}). Three pieces of that picture have no direct
Mathlib counterpart and must be built here, each in its own section:

\begin{enumerate}
  \item \textbf{Tensor powers of a sheaf of modules}
    (\cref{sec:sgr_tensor_powers}). Mathlib equips the category
    \(\PshMod\) of \emph{presheaves} of modules over a presheaf of commutative
    rings with a symmetric monoidal structure
    (\cref{lem:presheafModule_monoidal_mathlib}) and provides the
    sheafification functor \(\PshMod \to \ShMod\)
    (\cref{lem:presheafModule_sheafification_mathlib}); but the category
    \(\ShMod\) of \emph{sheaves} of modules carries \emph{no} monoidal
    structure. The tensor product of sheaves of modules
    \(\mathcal{F} \otimes_{\mathcal{O}_X} \mathcal{G}\) must therefore be built
    as the sheafification of the pointwise tensor presheaf, and the iterated
    tensor power \(\mathcal{L}^{\otimes m}\) recursively from it.
  \item \textbf{Lax-monoidal global sections} (\cref{sec:sgr_lax_sections}).
    The global-sections functor \(\Gamma\) is only \emph{lax} monoidal: from the
    pointwise-strict monoidality of the presheaf tensor product and the
    sheafification unit one obtains a natural multiplication
    \(\Gamma(\mathcal{F}) \otimes_{\Gamma(\mathcal{O}_X)} \Gamma(\mathcal{G})
    \to \Gamma(\mathcal{F} \otimes_{\mathcal{O}_X} \mathcal{G})\), together with
    the associativity and unit coherence those products must satisfy.
  \item \textbf{Graded-ring / graded-module assembly}
    (\cref{sec:sgr_graded_assembly}). The section multiplications and the
    tensor-power comparison isomorphisms assemble, through Mathlib's
    external-direct-sum graded API (\cref{lem:directSum_gcommSemiring_mathlib},
    \cref{lem:directSum_gmodule_mathlib}), into the graded \emph{ring} structure
    on \(\bigoplus_{m \ge 0} \Gamma(X_s, L_s^{\otimes m})\) and the graded
    \emph{module} structure on
    \(\bigoplus_{m \ge 0} \Gamma(X_s, \mathcal{F}_s \otimes L_s^{\otimes m})\).
\end{enumerate}

\section{Tensor powers of a sheaf of modules}
\label{sec:sgr_tensor_powers}

We work over a fixed scheme \(X\) (in the application \(X = X_s\), a fibre over a
field \(\kappa(s)\)). Its category of sheaves of modules is
\(X.\Modules = \ShMod(\mathcal{O}_X)\), with global-sections notation
\(\Gamma(X, \mathcal{M}) = \mathcal{M}(X)\) for the value of the underlying
presheaf on the top open. The structure sheaf \(\mathcal{O}_X\) is a sheaf of
commutative rings, so its underlying presheaf factors through the category of
commutative rings; this is the hypothesis under which Mathlib's presheaf-level
monoidal structure applies.

\begin{lemma}[Monoidal structure on presheaves of modules]
  \label{lem:presheafModule_monoidal_mathlib}
  \lean{PresheafOfModules.monoidalCategory}
  \mathlibok
  \textit{Provided by Mathlib
  (\texttt{Mathlib.Algebra.Category.ModuleCat.Presheaf.Monoidal}).}
  Let \(R : \mathcal{C}^{\mathrm{op}} \to \mathrm{CommRingCat}\) be a presheaf of
  commutative rings. Then the category \(\PshMod(R)\) of presheaves of modules
  over \(R\) carries a symmetric monoidal structure whose tensor product is
  computed objectwise: for presheaves of modules \(\mathcal{M}_1, \mathcal{M}_2\)
  and an object \(U\),
  \[
    (\mathcal{M}_1 \otimes \mathcal{M}_2)(U)
      \;=\; \mathcal{M}_1(U) \otimes_{R(U)} \mathcal{M}_2(U),
  \]
  with restriction maps the tensor products of the restriction maps. The unit is
  the presheaf \(R\) viewed as a module over itself, and the associator,
  unitors, and braiding are induced objectwise from those of
  \(\mathrm{ModuleCat}\). In particular evaluation at a fixed object \(U\) is a
  \emph{strict} monoidal functor
  \(\PshMod(R) \to \mathrm{ModuleCat}(R(U))\).
\end{lemma}

\begin{lemma}[Sheafification of presheaves of modules]
  \label{lem:presheafModule_sheafification_mathlib}
  \lean{PresheafOfModules.sheafification}
  \mathlibok
  \textit{Provided by Mathlib
  (\texttt{Mathlib.Algebra.Category.ModuleCat.Presheaf.Sheafification}).}
  Let \(R\) be a sheaf of (commutative) rings on a site, with underlying
  presheaf \(R_0\). There is a sheafification functor
  \[
    (-)^{\#} \;:\; \PshMod(R_0) \longrightarrow \ShMod(R),
  \]
  left adjoint to the forgetful functor, together with the unit
  \(\eta_{\mathcal{M}_0} : \mathcal{M}_0 \to (\mathcal{M}_0)^{\#}\) of the
  adjunction. On underlying sheaves of abelian groups it agrees with the
  ordinary sheafification, and it sends a presheaf of \(R_0\)-modules to the
  associated sheaf of \(R\)-modules.
\end{lemma}

\begin{lemma}[The structure sheaf as the unit module]
  \label{lem:moduleUnit_mathlib}
  \lean{SheafOfModules.unit}
  \mathlibok
  \textit{Provided by Mathlib
  (\texttt{Mathlib.Algebra.Category.ModuleCat.Sheaf}).}
  For a sheaf of rings \(R\) (in the application \(R = \mathcal{O}_X\)) the
  structure sheaf, viewed as a module over itself, is a distinguished object
  \(\mathbf{1} = \mathcal{O}_X \in \ShMod(\mathcal{O}_X)\). It is the value at
  the zeroth tensor power, \(\mathcal{L}^{\otimes 0} = \mathcal{O}_X\), and its
  global sections \(\Gamma(X, \mathcal{O}_X)\) form the degree-zero component of
  the section graded ring.
\end{lemma}


\begin{definition}\leanok
[Unit module of a scheme]
  \label{def:unitModule}
  \lean{AlgebraicGeometry.Scheme.Modules.unitModule}
  \uses{lem:moduleUnit_mathlib}
  For a scheme \(X\), the \emph{unit module} \(\mathbf{1}_X \in X.\Modules\) is the
  structure sheaf viewed as a module over itself, i.e.\ Mathlib's
  \(\mathrm{SheafOfModules.unit}\) (\cref{lem:moduleUnit_mathlib}) for
  \(\mathcal{O}_X\). It is the value of the zeroth tensor power,
  \(\mathcal{L}^{\otimes 0} = \mathbf{1}_X\) (\cref{def:sheafTensorPow}).
\end{definition}

\begin{definition}\leanok
[Tensor product of sheaves of modules]
  \label{def:sheafTensorObj}
  \lean{AlgebraicGeometry.Scheme.Modules.tensorObj}
  \uses{lem:presheafModule_monoidal_mathlib,
    lem:presheafModule_sheafification_mathlib}
  % SOURCE: [Stacks Project], "Sheaves of Modules", \S "Tensor product" (Tag
  % 01CA), prose definition (read from references/stacks-modules.tex,
  % L2282-L2297).
  % SOURCE QUOTE:
  % "Let $(X, \mathcal{O}_X)$ be a ringed space and let
  %  $\mathcal{F}$ and $\mathcal{G}$ be $\mathcal{O}_X$-modules.
  %  We define first the {\it tensor product presheaf}
  %  $$
  %  \mathcal{F} \otimes_{p, \mathcal{O}_X} \mathcal{G}
  %  $$
  %  as the rule which assigns to $U \subset X$ open the
  %  $\mathcal{O}_X(U)$-module
  %  $\mathcal{F}(U) \otimes_{\mathcal{O}_X(U)} \mathcal{G}(U)$.
  %  Having defined this we define the {\it tensor product sheaf}
  %  as the sheafification of the above:
  %  $$
  %  \mathcal{F} \otimes_{\mathcal{O}_X} \mathcal{G}
  %  =
  %  (\mathcal{F} \otimes_{p, \mathcal{O}_X} \mathcal{G})^\#
  %  $$"
  \textit{Source: [Stacks Project], "Sheaves of Modules", Tag 01CA (tensor
  product).}
  Let \(\mathcal{F}, \mathcal{G}\) be sheaves of \(\mathcal{O}_X\)-modules. Their
  \emph{tensor product} is the sheaf of \(\mathcal{O}_X\)-modules
  \[
    \mathcal{F} \otimes_{\mathcal{O}_X} \mathcal{G}
      \;=\; \bigl(\mathcal{F} \otimes_{p,\mathcal{O}_X} \mathcal{G}\bigr)^{\#},
  \]
  the sheafification (\cref{lem:presheafModule_sheafification_mathlib}) of the
  tensor product presheaf \(\mathcal{F} \otimes_{p,\mathcal{O}_X} \mathcal{G}\),
  whose value on an open \(U\) is the \(\mathcal{O}_X(U)\)-module
  \(\mathcal{F}(U) \otimes_{\mathcal{O}_X(U)} \mathcal{G}(U)\). Concretely, the
  tensor product presheaf is the Mathlib presheaf-level monoidal product
  (\cref{lem:presheafModule_monoidal_mathlib}) of the underlying presheaves of
  modules of \(\mathcal{F}\) and \(\mathcal{G}\), and
  \(\mathcal{F} \otimes_{\mathcal{O}_X} \mathcal{G}\) is its sheafification. This
  construction is symmetric and associative because the presheaf tensor product
  is, and sheafification preserves the coherence isomorphisms.
\end{definition}

\begin{definition}\leanok
[Tensor power of a line bundle]
  \label{def:sheafTensorPow}
  \lean{AlgebraicGeometry.Scheme.Modules.tensorPow}
  \uses{def:sheafTensorObj, def:unitModule, lem:moduleUnit_mathlib}
  % SOURCE: [Stacks Project], "Sheaves of Modules", \S "Invertible modules",
  % `definition-powers' (Tag 01CU) (read from references/stacks-modules.tex,
  % L4220-L4228).
  % SOURCE QUOTE:
  % "Let $(X, \mathcal{O}_X)$ be a ringed space. Given an invertible sheaf
  %  $\mathcal{L}$ on $X$ and $n \in \mathbf{Z}$ we define the
  %  $n$th {\it tensor power} $\mathcal{L}^{\otimes n}$ of $\mathcal{L}$
  %  as the image of $\mathcal{O}_X$ under applying the equivalence
  %  $\mathcal{F} \mapsto\mathcal{F} \otimes_{\mathcal{O}_X} \mathcal{L}$
  %  exactly $n$ times."
  \textit{Source: [Stacks Project], "Sheaves of Modules", Tag 01CU (tensor
  powers).}
  Let \(\mathcal{L}\) be a sheaf of \(\mathcal{O}_X\)-modules (in the application
  an invertible sheaf / line bundle \(L_s\)). For \(m \in \mathbb{N}\) the
  \emph{\(m\)th tensor power} \(\mathcal{L}^{\otimes m}\) is defined by recursion
  on \(m\):
  \[
    \mathcal{L}^{\otimes 0} = \mathcal{O}_X
      \quad(\text{the unit module, \cref{lem:moduleUnit_mathlib}}),
    \qquad
    \mathcal{L}^{\otimes (m+1)}
      = \mathcal{L}^{\otimes m} \otimes_{\mathcal{O}_X} \mathcal{L}
        \quad(\cref{def:sheafTensorObj}).
  \]
  We restrict to non-negative powers, which is all the section graded ring uses;
  the Stacks definition extends \(\mathcal{L}^{\otimes n}\) to all
  \(n \in \mathbf{Z}\) using invertibility, but negative powers are not needed
  here. For \(\mathcal{L}\) invertible each \(\mathcal{L}^{\otimes m}\) is again
  invertible.
\end{definition}


\begin{definition}\leanok
[Twist of a coherent sheaf by a tensor power]
  \label{def:sheafModuleTwist}
  \lean{AlgebraicGeometry.Scheme.Modules.moduleTensorPow}
  \uses{def:sheafTensorObj, def:sheafTensorPow}
  Let \(\mathcal{F}\) be a sheaf of \(\mathcal{O}_X\)-modules (in the application
  a coherent sheaf \(\mathcal{F}_s\)) and \(\mathcal{L}\) a line bundle. For
  \(m \in \mathbb{N}\) the \emph{\(m\)-twist} of \(\mathcal{F}\) by
  \(\mathcal{L}\) is
  \[
    \mathcal{F}(m) \;:=\;
      \mathcal{F} \otimes_{\mathcal{O}_X} \mathcal{L}^{\otimes m}
      \qquad(\cref{def:sheafTensorObj}, \cref{def:sheafTensorPow}),
  \]
  the degree-\(m\) carrier of the section graded module. By construction
  \(\mathcal{F}(0) = \mathcal{F} \otimes_{\mathcal{O}_X} \mathcal{O}_X \cong
  \mathcal{F}\).
\end{definition}


The next four isomorphisms are the parts of the (would-be) symmetric monoidal
structure on \(X.\Modules\) that descend through sheafification from
\cref{lem:presheafModule_monoidal_mathlib} using only \emph{functoriality} of the
sheafification functor (\cref{def:schemeModuleSheafification}) and, for the
unitors, the reflective counit isomorphism --- no strong monoidality of
sheafification is required. They are the structural inputs of the tensor-power
comparison \cref{lem:sheafTensorPow_add}.


The strong-monoidality comparison \cref{lem:isIso_sheafification_whiskerRight_unit}
below is assembled from three reduction lemmas --- a localization criterion, the
local-isomorphism property of the unit, and their corollary --- together with the
coequalizer presentation of the relative module-presheaf tensor product, the single
Mathlib-absent brick on which the comparison rests.


\begin{definition}

\leanok
[Scheme-level module sheafification functor]
  \label{def:schemeModuleSheafification}
  \lean{AlgebraicGeometry.Scheme.Modules.sheafification}
  \uses{lem:presheafModule_sheafification_mathlib}
  For a scheme \(X\), the \emph{module sheafification functor}
  \[
    (-)^{\#} \;:\; X.\PshMod \longrightarrow X.\Modules
  \]
  is Mathlib's presheaf-of-modules sheafification
  (\cref{lem:presheafModule_sheafification_mathlib}) for the identity morphism of
  the underlying presheaf of (commutative) rings of \(\mathcal{O}_X\), which is
  locally bijective. It is left adjoint to the forgetful functor
  \(X.\Modules \to X.\PshMod\), with adjunction unit
  \(\eta_{\mathcal{M}_0} : \mathcal{M}_0 \to (\mathcal{M}_0)^{\#}\) and an
  invertible counit (the forgetful functor is fully faithful). Being a functor it
  carries isomorphisms to isomorphisms, which is precisely what lets the
  presheaf-level coherence isomorphisms of
  \cref{lem:presheafModule_monoidal_mathlib} descend to \(X.\Modules\).
\end{definition}

\begin{definition}

\leanok
[Left unitor of the sheaf tensor product]
  \label{def:tensorObjUnitIso}
  \lean{AlgebraicGeometry.Scheme.Modules.tensorObjUnitIso}
  \uses{def:sheafTensorObj, def:unitModule, def:schemeModuleSheafification,
    lem:presheafModule_monoidal_mathlib}
  For a sheaf of \(\mathcal{O}_X\)-modules \(G\), the \emph{left unitor}
  \[
    \mathbf{1}_X \otimes_{\mathcal{O}_X} G \;\xrightarrow{\ \sim\ }\; G
  \]
  is obtained by descending the presheaf-level left unitor of
  \cref{lem:presheafModule_monoidal_mathlib} through the sheafification functor
  (\cref{def:schemeModuleSheafification}) and composing with the reflective counit
  isomorphism \((\mathcal{M}_0)^{\#} \cong \mathcal{M}\) of the sheafification
  adjunction. It is the base case \(m = 0\) of the tensor-power comparison
  \cref{lem:sheafTensorPow_add}, and uses no monoidal structure on
  \(X.\Modules\) --- only functoriality of sheafification and the invertible
  counit.
\end{definition}

\begin{definition}

\leanok
[Braiding of the sheaf tensor product]
  \label{def:tensorBraiding}
  \lean{AlgebraicGeometry.Scheme.Modules.tensorBraiding}
  \uses{def:sheafTensorObj, def:schemeModuleSheafification,
    lem:presheafModule_monoidal_mathlib}
  For sheaves of \(\mathcal{O}_X\)-modules \(F, G\), the \emph{braiding}
  \[
    F \otimes_{\mathcal{O}_X} G \;\xrightarrow{\ \sim\ }\;
      G \otimes_{\mathcal{O}_X} F
  \]
  is the image under the sheafification functor
  (\cref{def:schemeModuleSheafification}) of the symmetric braiding of the
  presheaf symmetric monoidal structure (\cref{lem:presheafModule_monoidal_mathlib})
  on the underlying presheaves of modules. As pure sheafification-functoriality of
  an isomorphism it is again an isomorphism, requiring no monoidal structure on
  \(X.\Modules\). It is the symmetry used in the inductive step of
  \cref{lem:sheafTensorPow_add}.
\end{definition}

\begin{lemma}

\leanok
[Sheafification inverts the whiskered localization unit]
  \label{lem:isIso_sheafification_whiskerRight_unit}
  \lean{AlgebraicGeometry.Scheme.Modules.isIso_sheafification_whiskerRight_unit}
  \uses{def:schemeModuleSheafification, lem:presheafModule_monoidal_mathlib,
    lem:presheafModule_sheafification_mathlib}
  Let \(P, Q\) be presheaves of \(\mathcal{O}_X\)-modules and let
  \(\eta_P : P \to P^{\#}\) be the unit of the sheafification adjunction
  (\cref{def:schemeModuleSheafification},
  \cref{lem:presheafModule_sheafification_mathlib}). Then the sheafification of the
  right-whiskered unit
  \[
    (\eta_P \rhd Q)^{\#} :
      (P \otimes_{p,\mathcal{O}_X} Q)^{\#} \longrightarrow
      (P^{\#} \otimes_{p,\mathcal{O}_X} Q)^{\#},
  \]
  where \(\otimes_{p,\mathcal{O}_X}\) is the presheaf monoidal product
  (\cref{lem:presheafModule_monoidal_mathlib}), is an isomorphism of sheaves of
  modules. This is the strong-monoidality comparison of the module sheafification
  functor on a whiskered unit, the single Mathlib-absent brick on which the
  sheaf-level associator \cref{cor:sheafTensorObjAssoc} and the tensor-power
  comparison \cref{lem:sheafTensorPow_add} rest.
\end{lemma}

\begin{proof}
  \uses{def:schemeModuleSheafification, lem:presheafModule_sheafification_mathlib,
    lem:presheafModule_monoidal_mathlib}
    \leanok
  Module sheafification is the localization functor at the class \(W\) of
  morphisms of presheaves of modules whose underlying morphism of abelian-group
  presheaves is a local isomorphism for the opens topology of \(X\). Hence, by the
  localization criterion, \((\eta_P \rhd Q)^{\#}\) is an isomorphism if and only if
  the underlying abelian-presheaf morphism of \(\eta_P \rhd Q\) is a local
  isomorphism. The unit \(\eta_P\) is, on underlying abelian presheaves, the
  ordinary sheafification unit \(\mathrm{toSheafify}\), which is a local
  isomorphism. Right-whiskering by \(Q\) preserves this: the relative
  module-presheaf tensor product is presented as a coequalizer, and tensoring with
  the fixed presheaf \(Q\) carries the local isomorphism \(\eta_P\) to a local
  isomorphism. Feeding this through the localization criterion gives the claim.
\end{proof}

\subsection{Monoidal structure by localization transport}
\label{subsec:sgr_monoidal_localization}

The preceding isomorphisms (unitors, braiding, whiskered-unit comparison) are
exactly the data Mathlib needs to \emph{transport} the entire symmetric monoidal
structure from presheaves of modules onto \(X.\Modules\) along the sheafification
functor, rather than reassembling associator, pentagon, triangle and hexagon by
hand. The mechanism is \emph{monoidal localization}: a localization functor whose
inverted class of morphisms is compatible with the tensor product carries the
monoidal structure of its source to its target, with all coherence laws inherited.
The two Mathlib anchors below supply the general transport theorem and the fact
that module sheafification is such a localization; the two project blocks assemble
the compatibility hypothesis and read off the resulting structure.

\begin{lemma}[Monoidal localization transport]
  \label{lem:monoidalLocalization_mathlib}
  \lean{CategoryTheory.Localization.LocalizedMonoidal}
  \mathlibok
  \textit{Provided by Mathlib
  (\texttt{Mathlib.CategoryTheory.Localization.Monoidal.Basic} and
  \texttt{\ldots.Localization.Monoidal.Braided}).}
  % SOURCE: [Mathlib], \texttt{Mathlib.CategoryTheory.Localization.Monoidal.Basic}
  %   (Mathlib/CategoryTheory/Localization/Monoidal/Basic.lean) and
  %   \texttt{Mathlib.CategoryTheory.Localization.Monoidal.Braided}
  %   (Mathlib/CategoryTheory/Localization/Monoidal/Braided.lean).
  % SOURCE QUOTE:
  % class MorphismProperty.IsMonoidal : Prop extends W.IsMultiplicative where
  %   whiskerLeft (X : C) {Y₁ Y₂ : C} (g : Y₁ ⟶ Y₂) (hg : W g) : W (X ◁ g)
  %   whiskerRight {X₁ X₂ : C} (f : X₁ ⟶ X₂) (hf : W f) (Y : C) : W (f ▷ Y)
  %
  % def LocalizedMonoidal (L : C ⥤ D) (W : MorphismProperty C)
  %     [W.IsMonoidal] [L.IsLocalization W] {unit : D} (_ : L.obj (𝟙_ C) ≅ unit) := D
  %
  % noncomputable def μ (X Y : C) : (L').obj X ⊗ (L').obj Y ≅ (L').obj (X ⊗ Y) :=
  %   ((tensorBifunctorIso L W ε).app X).app Y
  Let \(\mathcal{C}\) be a monoidal category and let \(W\) be a class of morphisms
  of \(\mathcal{C}\) satisfying \(\mathrm{MorphismProperty.IsMonoidal}\): \(W\) is
  multiplicative and stable under left- and right-whiskering by arbitrary objects.
  Let \(L : \mathcal{C} \to \mathcal{D}\) be a
  localization functor for \(W\). Then \(\mathcal{D}\) acquires a full monoidal
  category structure --- tensor product, associator and unitors --- for which the
  pentagon and triangle coherence identities hold, and \(L\) is strong monoidal,
  with comparison isomorphisms
  \[
    \mu_{X,Y} : L(X) \otimes L(Y) \;\xrightarrow{\ \sim\ }\; L(X \otimes Y).
  \]
  All of this structure is inherited from \(\mathcal{C}\) through \(L\). If moreover
  \(\mathcal{C}\) is braided (resp.\ symmetric), then \(\mathcal{D}\) inherits a
  braided (resp.\ symmetric) structure for which the hexagon identities hold and
  \(L\) is braided monoidal. The sheafification of presheaves over a site is the
  archetypal instance of this transport.
\end{lemma}

\begin{lemma}[Module sheafification is a monoidal localization input]
  \label{lem:moduleSheafification_isLocalization_mathlib}
  \lean{PresheafOfModules.sheafification}
  \mathlibok
  \textit{Provided by Mathlib
  (\texttt{Mathlib.Algebra.Category.ModuleCat.Sheaf.Localization} and
  \texttt{\ldots.ModuleCat.Presheaf.Monoidal}).}
  % SOURCE: [Mathlib], \texttt{Mathlib.Algebra.Category.ModuleCat.Sheaf.Localization}
  %   (Mathlib/Algebra/Category/ModuleCat/Sheaf/Localization.lean).
  % SOURCE QUOTE:
  % instance : (sheafification.{v} α).IsLocalization (J.W.inverseImage (toPresheaf R₀)) := by
  %   rw [inverseImage_W_toPresheaf_eq_inverseImage_isomorphisms α]
  %   exact (sheafificationAdjunction.{v} α).isLocalization
  The sheafification functor for presheaves of modules
  (\cref{lem:presheafModule_sheafification_mathlib}) is a \emph{localization
  functor} for the morphism class
  \[
    W' \;=\; J.W.\mathrm{inverseImage}\,(\mathrm{toPresheaf}\,R_0),
  \]
  the class of morphisms of presheaves of \(R_0\)-modules whose underlying
  morphism of abelian-group presheaves lies in the local-isomorphism class
  \(J.W\) of the topology --- equivalently, the morphisms that sheafification
  sends to isomorphisms. Moreover the source category \(\PshMod(R_0)\) is
  symmetric monoidal (\cref{lem:presheafModule_monoidal_mathlib}). Thus
  sheafification supplies precisely a localization functor \(L\) on a symmetric
  monoidal source, the two standing inputs of the transport theorem
  \cref{lem:monoidalLocalization_mathlib}; the only remaining hypothesis is that
  \(W'\) be compatible with the tensor product.
\end{lemma}

\begin{definition}
[Tensor-compatibility of the sheafification localization class]
  \label{def:sheafModule_W_isMonoidal}
  \lean{AlgebraicGeometry.Scheme.Modules.W_isMonoidal}
  \uses{lem:isIso_sheafification_whiskerRight_unit,
    lem:presheafModule_monoidal_mathlib,
    lem:moduleSheafification_isLocalization_mathlib}
  For a scheme \(X\), write \(R_0\) for the underlying presheaf of the structure
  sheaf \(\mathcal{O}_X\) and
  \(W' = J.W.\mathrm{inverseImage}\,(\mathrm{toPresheaf}\,R_0)\) for the
  sheafification localization class of
  \cref{lem:moduleSheafification_isLocalization_mathlib}. Then \(W'\) satisfies
  \(\mathrm{MorphismProperty.IsMonoidal}\): it is multiplicative and stable under
  left- and right-whiskering. This is the compatibility hypothesis that, together
  with \cref{lem:moduleSheafification_isLocalization_mathlib}, feeds the transport
  theorem \cref{lem:monoidalLocalization_mathlib}.
\end{definition}

\begin{proof}
  \uses{lem:isIso_sheafification_whiskerRight_unit,
    lem:presheafModule_monoidal_mathlib,
    lem:moduleSheafification_isLocalization_mathlib}
  Multiplicativity of \(W'\) is free: \(W'\) is the inverse image of a class along
  a functor, and the inverse image of a multiplicative class is multiplicative.
  For the right-whiskering field, a morphism \(f\) lies in \(W'\) precisely when
  sheafification inverts it, and the general whisker comparison underlying
  \cref{lem:isIso_sheafification_whiskerRight_unit} --- stated there for the unit,
  but established for an arbitrary morphism inverted by sheafification --- shows
  that \(f \rhd Q\) is again inverted by sheafification, hence lies in \(W'\); the
  bridge between ``lies in \(W'\)'' and ``sheafification inverts'' is the
  localization criterion for the sheafification functor. The left-whiskering field
  follows by conjugating the right-whiskering field with the symmetric braiding of
  the presheaf monoidal structure (\cref{lem:presheafModule_monoidal_mathlib}):
  since \(Q \lhd f\) is carried to \(f \rhd Q\) by the braiding isomorphism on
  both ends, and membership in \(W'\) is invariant under isomorphism of arrows,
  stability under left-whiskering reduces to stability under right-whiskering.
\end{proof}

\begin{definition}
[Symmetric monoidal structure on sheaves of modules by transport]
  \label{def:sheafModule_monoidalStructure}
  \lean{AlgebraicGeometry.Scheme.Modules.monoidalCategory}
  \uses{def:sheafModule_W_isMonoidal, lem:monoidalLocalization_mathlib}
  For a scheme \(X\), the category \(X.\Modules\) of sheaves of
  \(\mathcal{O}_X\)-modules carries a symmetric monoidal structure obtained by
  transporting the symmetric monoidal structure of \(X.\PshMod\) along the
  sheafification functor. Concretely, one instantiates the monoidal localization
  transport \cref{lem:monoidalLocalization_mathlib} with source the symmetric
  monoidal category \(X.\PshMod\), localization functor the module sheafification
  \cref{lem:moduleSheafification_isLocalization_mathlib}, and the
  tensor-compatibility of its inverted class \cref{def:sheafModule_W_isMonoidal};
  the unit is chosen to be the unit module \(\mathbf{1}_X\) (\cref{def:unitModule}).
  The resulting tensor product agrees objectwise with the sheaf tensor product
  \cref{def:sheafTensorObj}, via the strong-monoidal comparison
  \(\mu\) of \cref{lem:monoidalLocalization_mathlib}. In particular the associator
  is the canonical (Mac Lane) associator, the pentagon and triangle laws hold, and
  --- the source being symmetric --- the braiding and its hexagon laws are
  inherited as well, giving \(X.\Modules\) the structure of a
  \emph{symmetric monoidal category}.
\end{definition}

\begin{proof}
  \uses{def:sheafModule_W_isMonoidal, lem:monoidalLocalization_mathlib}
  By \cref{lem:moduleSheafification_isLocalization_mathlib} the sheafification
  functor is a localization for \(W'\) on the symmetric monoidal source
  \(X.\PshMod\), and by \cref{def:sheafModule_W_isMonoidal} the class \(W'\)
  satisfies \(\mathrm{MorphismProperty.IsMonoidal}\). The hypotheses of the
  transport theorem \cref{lem:monoidalLocalization_mathlib} are therefore met, so
  \(X.\Modules\) acquires a monoidal structure with pentagon and triangle, and ---
  the source being symmetric --- a symmetric structure with the hexagon
  identities, all inherited through sheafification. The three coherence
  identities required downstream for the section-graded-ring assembly hence hold by
  inheritance rather than by separate induction.
\end{proof}

\begin{corollary}

\leanok
[Associator of the sheaf tensor product]
  \label{cor:sheafTensorObjAssoc}
  \lean{AlgebraicGeometry.Scheme.Modules.tensorObjAssoc}
  \uses{def:sheafTensorObj, def:schemeModuleSheafification,
    lem:isIso_sheafification_whiskerRight_unit,
    lem:presheafModule_monoidal_mathlib}
  For sheaves of \(\mathcal{O}_X\)-modules \(A, B, C\) there is a canonical
  associativity isomorphism
  \[
    (A \otimes_{\mathcal{O}_X} B) \otimes_{\mathcal{O}_X} C
      \;\xrightarrow{\ \sim\ }\;
      A \otimes_{\mathcal{O}_X} (B \otimes_{\mathcal{O}_X} C)
  \]
  of sheaves of \(\mathcal{O}_X\)-modules.
\end{corollary}

\begin{proof}
  \uses{def:sheafTensorObj, def:schemeModuleSheafification,
    lem:isIso_sheafification_whiskerRight_unit,
    lem:presheafModule_monoidal_mathlib}
    \leanok
  Each iterated sheaf tensor product carries inner sheafifications
  (\cref{def:sheafTensorObj}); writing \(a, b, c\) for the underlying presheaves of
  modules of \(A, B, C\),
  \[
    (A \otimes B) \otimes C =
      \bigl((a \otimes_{p,\mathcal{O}_X} b)^{\#}
        \otimes_{p,\mathcal{O}_X} c\bigr)^{\#},
    \qquad
    A \otimes (B \otimes C) =
      \bigl(a \otimes_{p,\mathcal{O}_X}
        (b \otimes_{p,\mathcal{O}_X} c)^{\#}\bigr)^{\#}.
  \]
  The associator is a three-segment composite. First, the inverse of the
  whiskered-unit isomorphism \cref{lem:isIso_sheafification_whiskerRight_unit}
  clears the inner sheafification on the left factor, identifying
  \((A \otimes B) \otimes C\) with the sheafification of the triple presheaf tensor
  \((a \otimes_{p,\mathcal{O}_X} b) \otimes_{p,\mathcal{O}_X} c\). Next, the
  presheaf-level associator (\cref{lem:presheafModule_monoidal_mathlib}) descends
  through sheafification (\cref{def:schemeModuleSheafification}) to reassociate to
  the sheafification of \(a \otimes_{p,\mathcal{O}_X}
  (b \otimes_{p,\mathcal{O}_X} c)\). Finally, a second application of
  \cref{lem:isIso_sheafification_whiskerRight_unit} --- conjugated by the presheaf
  braiding so as to act on the right factor --- reinstates the inner sheafification,
  landing in \(A \otimes (B \otimes C)\). Each segment is an isomorphism, hence so
  is the composite.
\end{proof}


\begin{lemma}\leanok
[Tensor-power comparison isomorphism]
  \label{lem:sheafTensorPow_add}
  \lean{AlgebraicGeometry.Scheme.Modules.tensorPowAdd}
  \uses{def:sheafTensorPow, def:sheafTensorObj}
  % SOURCE: [Stacks Project], "Sheaves of Modules", \S "Invertible modules"
  % (Tag 01CU, following `definition-powers') (read from
  % references/stacks-modules.tex, L4249-L4254).
  % SOURCE QUOTE:
  % "With this definition we have canonical isomorphisms
  %  $\mathcal{L}^{\otimes n} \otimes_{\mathcal{O}_X}
  %  \mathcal{L}^{\otimes m} \to
  %  \mathcal{L}^{\otimes n + m}$, and these isomorphisms
  %  satisfy a commutativity and an associativity constraint
  %  (formulation omitted)."
  \textit{Source: [Stacks Project], "Sheaves of Modules", Tag 01CU (canonical
  power isomorphisms).}
  For a line bundle \(\mathcal{L}\) and \(m, m' \in \mathbb{N}\) there is a
  canonical isomorphism of sheaves of \(\mathcal{O}_X\)-modules
  \[
    \mu_{m,m'} :
    \mathcal{L}^{\otimes m} \otimes_{\mathcal{O}_X} \mathcal{L}^{\otimes m'}
    \;\xrightarrow{\ \sim\ }\; \mathcal{L}^{\otimes (m+m')},
  \]
  natural in \(\mathcal{L}\), and these isomorphisms satisfy a commutativity
  constraint (\(\mu_{m,m'}\) agrees with \(\mu_{m',m}\) after the braiding) and
  an associativity constraint (the two ways of combining
  \(\mathcal{L}^{\otimes m} \otimes \mathcal{L}^{\otimes m'} \otimes
  \mathcal{L}^{\otimes m''}\) into \(\mathcal{L}^{\otimes (m+m'+m'')}\) agree).
\end{lemma}

\begin{proof}
  \uses{def:sheafTensorPow, def:sheafTensorObj, cor:sheafTensorObjAssoc,
    def:tensorObjUnitIso, def:tensorBraiding,
    lem:isIso_sheafification_whiskerRight_unit,
    lem:presheafModule_monoidal_mathlib,
    lem:presheafModule_sheafification_mathlib}
    \leanok
  We do \emph{not} build a full monoidal structure on the whole category
  \(\ShMod(\mathcal{O}_X)\) --- none is constructed here --- but only assemble
  the comparison \(\mu_{m,m'}\) from those structural isomorphisms of
  \cref{def:sheafTensorObj} that the induction actually uses: the associator,
  the braiding, and the left unitor.

  \textbf{The structural maps.} The braiding and the unitors descend to
  \(\ShMod(\mathcal{O}_X)\) directly: each is the sheafification of the
  corresponding isomorphism of the presheaf symmetric monoidal structure
  (\cref{lem:presheafModule_monoidal_mathlib}), and sheafification, being a
  functor, carries isomorphisms to isomorphisms. The \emph{associator} is not
  the sheafification of a single presheaf isomorphism, because the iterated
  sheaf tensor products
  \((\mathcal{A} \otimes \mathcal{B}) \otimes \mathcal{C}\) and
  \(\mathcal{A} \otimes (\mathcal{B} \otimes \mathcal{C})\) each carry an inner
  sheafification (\cref{def:sheafTensorObj}); it is supplied instead by
  \cref{cor:sheafTensorObjAssoc}, whose construction clears those inner
  sheafifications using the strong-monoidality comparison
  \cref{lem:isIso_sheafification_whiskerRight_unit}. These are the only
  isomorphisms of sheaves of modules the construction invokes; right-exactness
  of the tensor product is never used.

  \textbf{Construction of \(\mu_{m,m'}\).} Fix \(m'\) and induct on \(m\). For
  \(m = 0\),
  \[
    \mathcal{L}^{\otimes 0} \otimes \mathcal{L}^{\otimes m'}
      = \mathcal{O}_X \otimes \mathcal{L}^{\otimes m'}
      \xrightarrow{\ \sim\ } \mathcal{L}^{\otimes m'}
  \]
  is the left unitor of \cref{def:sheafTensorObj}. For the inductive step the
  recursive definition
  \(\mathcal{L}^{\otimes (m+1)} = \mathcal{L}^{\otimes m} \otimes \mathcal{L}\)
  (\cref{def:sheafTensorPow}) gives
  \[
    \mathcal{L}^{\otimes (m+1)} \otimes \mathcal{L}^{\otimes m'}
    = (\mathcal{L}^{\otimes m} \otimes \mathcal{L}) \otimes
      \mathcal{L}^{\otimes m'}
    \xrightarrow{\ \alpha\ } \mathcal{L}^{\otimes m} \otimes
      (\mathcal{L} \otimes \mathcal{L}^{\otimes m'})
    \xrightarrow{\ \mathrm{id} \otimes \beta\ } \mathcal{L}^{\otimes m} \otimes
      (\mathcal{L}^{\otimes m'} \otimes \mathcal{L}),
  \]
  using the associator \(\alpha\) (\cref{cor:sheafTensorObjAssoc}) and then the
  braiding \(\beta\) of \cref{def:sheafTensorObj} whiskered into the left
  factor. The inverse associator \(\alpha^{-1}\)
  (\cref{cor:sheafTensorObjAssoc}) and the inductive comparison \(\mu_{m,m'}\)
  whiskered by \(\mathcal{L}\) on the right then carry this to
  \[
    (\mathcal{L}^{\otimes m} \otimes \mathcal{L}^{\otimes m'}) \otimes
      \mathcal{L}
    \xrightarrow{\ \mu_{m,m'} \rhd \mathcal{L}\ }
    \mathcal{L}^{\otimes (m + m')} \otimes \mathcal{L}
    = \mathcal{L}^{\otimes ((m + m') + 1)},
  \]
  the final equality being the recursive definition. The reindexing
  \((m + m') + 1 = (m + 1) + m'\) identifies the target with
  \(\mathcal{L}^{\otimes ((m+1) + m')}\), and the composite is
  \(\mu_{m+1,m'}\), an isomorphism as a composite of isomorphisms. Naturality in
  \(\mathcal{L}\) holds because each constituent comparison is natural.

  \textbf{Commutativity and associativity constraints.} The braiding and
  unitors are the sheafifications of the symmetric-monoidal coherence data of
  \(\PshMod(\mathcal{O}_X)\) (\cref{lem:presheafModule_monoidal_mathlib}), and
  the associator is intertwined with the presheaf associator by the comparison
  isomorphisms of \cref{lem:isIso_sheafification_whiskerRight_unit} through the
  construction of \cref{cor:sheafTensorObjAssoc}. Both the commutativity square
  (\(\mu_{m,m'}\) versus \(\mu_{m',m}\) postcomposed with the braiding) and the
  associativity pentagon (the two bracketings of
  \(\mathcal{L}^{\otimes m} \otimes \mathcal{L}^{\otimes m'} \otimes
  \mathcal{L}^{\otimes m''} \to \mathcal{L}^{\otimes (m+m'+m'')}\)) are therefore
  images, under the sheafification functor together with these comparison
  isomorphisms, of the corresponding coherence diagrams of the presheaf
  symmetric monoidal structure, which commute by Mac\,Lane coherence. Since a
  functor preserves commuting diagrams and the comparison isomorphisms are
  natural, the constraints hold for the family \(\{\mu_{m,m'}\}\). We emphasise
  that this argument uses only the single strong-monoidality comparison
  \cref{lem:isIso_sheafification_whiskerRight_unit} and the resulting associator
  \cref{cor:sheafTensorObjAssoc}, not a full monoidal structure on
  \(\ShMod(\mathcal{O}_X)\).
\end{proof}

\section{Lax-monoidal global sections}
\label{sec:sgr_lax_sections}

The graded multiplication is supplied by the lax-monoidal structure of the
global-sections functor \(\Gamma(X, -) : \ShMod(\mathcal{O}_X) \to
\mathrm{ModuleCat}(\Gamma(X, \mathcal{O}_X))\). It is only \emph{lax} because
global sections do not commute with sheafification: a global section of a
sheafified tensor product is not in general a finite sum of elementary tensors of
global sections.

\begin{definition}\leanok
[Section multiplication]
  \label{def:sectionMul}
  \lean{AlgebraicGeometry.Scheme.Modules.sectionsMul}
  \uses{def:sheafTensorObj, lem:presheafModule_monoidal_mathlib,
    lem:presheafModule_sheafification_mathlib}
  Let \(\mathcal{F}, \mathcal{G}\) be sheaves of \(\mathcal{O}_X\)-modules. The
  \emph{section multiplication} is the \(\Gamma(X,\mathcal{O}_X)\)-bilinear map
  \[
    \Gamma(X, \mathcal{F}) \otimes_{\Gamma(X,\mathcal{O}_X)}
      \Gamma(X, \mathcal{G})
    \;\longrightarrow\;
    \Gamma(X, \mathcal{F} \otimes_{\mathcal{O}_X} \mathcal{G})
  \]
  defined as follows. A pair of global sections \((\sigma, \tau)\) determines,
  by the objectwise formula of \cref{lem:presheafModule_monoidal_mathlib} at the
  top open, the elementary tensor \(\sigma \otimes \tau \in
  (\mathcal{F} \otimes_{p,\mathcal{O}_X} \mathcal{G})(X)\), a global section of
  the tensor product presheaf. Composing with the global-sections component of
  the sheafification unit
  \(\eta : \mathcal{F} \otimes_{p,\mathcal{O}_X} \mathcal{G} \to
  (\mathcal{F} \otimes_{p,\mathcal{O}_X} \mathcal{G})^{\#} =
  \mathcal{F} \otimes_{\mathcal{O}_X} \mathcal{G}\)
  (\cref{lem:presheafModule_sheafification_mathlib},
  \cref{def:sheafTensorObj}) yields a global section
  \(\sigma \cdot \tau \in
  \Gamma(X, \mathcal{F} \otimes_{\mathcal{O}_X} \mathcal{G})\). Bilinearity over
  \(\Gamma(X,\mathcal{O}_X)\) is inherited from bilinearity of the elementary
  tensor and \(\Gamma(X,\mathcal{O}_X)\)-linearity of \(\eta\).
\end{definition}

The graded-monoid axioms below are not raw equalities between sections of two
different tensor powers --- on \(\mathbb{N}\) the indices \((i+j)+k\) and
\(i+(j+k)\), or \(0+n\) and \(n\), are equal but not definitionally so, so the two
sides a priori inhabit distinct section modules. Following
\texttt{Mathlib.LinearAlgebra.TensorPower.Basic}, we phrase each coherence as an
honest equality in a single section module, obtained by transporting one side
along an index-equality isomorphism, and provide a bridge that repackages such a
transported equality as an equality of dependent pairs.

\begin{definition}\leanok
[Section component in degree \(m\)]
  \label{def:sectionDeg}
  \lean{AlgebraicGeometry.Scheme.Modules.sectionDeg}
  \uses{def:sheafTensorPow}
  Let \(\mathcal{L}\) be a sheaf of \(\mathcal{O}_X\)-modules. The
  \emph{degree-\(m\) section component} is the \(\Gamma(X,\mathcal{O}_X)\)-module of
  global sections of the \(m\)th tensor power,
  \[
    (\mathrm{sectionDeg}\ \mathcal{L})_m
      \;:=\; \Gamma(X, \mathcal{L}^{\otimes m})
      \;=\; \mathcal{L}^{\otimes m}(X)
        \qquad(\cref{def:sheafTensorPow}).
  \]
  It inherits its additive-group and \(\Gamma(X,\mathcal{O}_X)\)-module structure
  from the underlying module object \(\mathcal{L}^{\otimes m}(X)\). These are the
  carriers \(m \mapsto (\mathrm{sectionDeg}\ \mathcal{L})_m\) of the section graded
  ring \(R(X_s, L_s)\) (\cref{lem:sectionGradedRing_gcommSemiring}) and the
  domains of the index transport \cref{def:sectionsCast}.
\end{definition}

\begin{definition}\leanok
[Index-equality transport of section components]
  \label{def:sectionsCast}
  \lean{AlgebraicGeometry.Scheme.Modules.sectionsCast}
  \uses{def:sheafTensorPow}
  Let \(\mathcal{L}\) be an invertible sheaf on \(X\) and let \(h : i = j\) be an
  equality of natural numbers. Applying the global-sections functor
  \(\Gamma(X,-)\) to the canonical isomorphism
  \(\mathcal{L}^{\otimes i} \xrightarrow{\sim} \mathcal{L}^{\otimes j}\) induced by
  \(h\) under the tensor-power functor (\cref{def:sheafTensorPow}) yields the
  \(\Gamma(X,\mathcal{O}_X)\)-linear isomorphism
  \[
    \mathrm{sectionsCast}(h) :
    \Gamma(X, \mathcal{L}^{\otimes i}) \;\xrightarrow{\ \sim\ }\;
    \Gamma(X, \mathcal{L}^{\otimes j}),
  \]
  the \emph{section-component transport} along \(h\). It is the section-level
  analogue of the tensor-power index cast \(\mathrm{TensorPower.cast}\) of
  \texttt{Mathlib.LinearAlgebra.TensorPower.Basic}, which transports
  \(\bigotimes^{i} M\) to \(\bigotimes^{j} M\) along \(i = j\). Because the
  isomorphism induced by the reflexive equality is the identity and
  \(\Gamma(X,-)\) preserves identities, the transport along a reflexive index
  equality is the identity map (the reflexivity simplification
  \(\mathrm{sectionsCast\_refl}\), \cref{lem:sectionsCast_refl}).
\end{definition}

\begin{lemma}\leanok
[Reflexivity of the section transport]
  \label{lem:sectionsCast_refl}
  \lean{AlgebraicGeometry.Scheme.Modules.sectionsCast_refl}
  \uses{def:sectionsCast}
  For any \(i\), the transport along the reflexive equality \(i = i\)
  (\cref{def:sectionsCast}) equals the identity automorphism of
  \(\Gamma(X, \mathcal{L}^{\otimes i})\). This is the section-level counterpart of
  \(\mathrm{TensorPower.cast\_refl}\), and serves as a simplification lemma
  collapsing the transport along a reflexive index equality.
\end{lemma}

\begin{proof}

\leanok
  The canonical isomorphism
  \(\mathcal{L}^{\otimes i} \cong \mathcal{L}^{\otimes i}\) induced by the
  reflexive equality is the identity, and the global-sections functor carries
  identities to identities.
\end{proof}

\begin{lemma}\leanok
[Cast-mediated equality in the graded monoid]
  \label{lem:gradedMonoid_eq_of_cast}
  \lean{AlgebraicGeometry.Scheme.Modules.gradedMonoid_eq_of_cast}
  \uses{def:sectionsCast, lem:sectionsCast_refl}
  Write the carrier of the graded monoid as the dependent sum
  \(\coprod_{n} \Gamma(X, \mathcal{L}^{\otimes n})\), whose elements are pairs
  \((n, s)\) with \(s \in \Gamma(X, \mathcal{L}^{\otimes n})\). Let
  \(a = (i, x)\) and \(b = (j, y)\) be two such pairs. If \(h : i = j\) and the
  transported component agrees, \(\mathrm{sectionsCast}(h)\,x = y\), then
  \(a = b\). That is, a component-level equality mediated by the index transport
  \cref{def:sectionsCast} repackages into a genuine equality of dependent pairs.
  This is the section-level analogue of the bridge
  \(\mathrm{gradedMonoid\_eq\_of\_cast}\) of
  \texttt{Mathlib.LinearAlgebra.TensorPower.Basic} (line~123 there).
\end{lemma}

\begin{proof}
  \uses{lem:sectionsCast_refl}
  \leanok
  Substituting \(j = i\) by means of \(h\) reduces the transport to the identity
  (\cref{lem:sectionsCast_refl}), so the hypothesis becomes \(x = y\); the two
  dependent pairs \((i, x)\) and \((i, y)\) then coincide.
\end{proof}

\begin{definition}\leanok
[Graded multiplication on section components]
  \label{def:sectionGradedGMul}
  \lean{AlgebraicGeometry.Scheme.Modules.instGMulNatSectionDeg}
  \uses{def:sectionDeg, def:sectionMul, lem:sheafTensorPow_add}
  The section components \cref{def:sectionDeg} carry a graded multiplication
  \[
    (\mathrm{sectionDeg}\ \mathcal{L})_i \times
      (\mathrm{sectionDeg}\ \mathcal{L})_j
    \;\longrightarrow\; (\mathrm{sectionDeg}\ \mathcal{L})_{i+j},
    \qquad (a, b) \longmapsto a \cdot b,
  \]
  landing in the index-sum component (a \(\mathrm{GradedMonoid.GMul}\) on
  \(m \mapsto (\mathrm{sectionDeg}\ \mathcal{L})_m\)). It is the section
  multiplication \(a \otimes b \mapsto a \cdot b\) (\cref{def:sectionMul}) on
  \(\mathcal{L}^{\otimes i}\) and \(\mathcal{L}^{\otimes j}\), followed by global
  sections of the tensor-power comparison isomorphism \(\mu_{i,j}\)
  (\cref{lem:sheafTensorPow_add}),
  \[
    \Gamma(X, \mathcal{L}^{\otimes i}) \otimes_{\Gamma(X,\mathcal{O}_X)}
      \Gamma(X, \mathcal{L}^{\otimes j})
    \xrightarrow{\ \cdot\ }
    \Gamma(X, \mathcal{L}^{\otimes i} \otimes_{\mathcal{O}_X}
      \mathcal{L}^{\otimes j})
    \xrightarrow{\ \Gamma(\mu_{i,j})\ }
    \Gamma(X, \mathcal{L}^{\otimes (i+j)}).
  \]
\end{definition}

\begin{definition}\leanok
[Graded unit in degree zero]
  \label{def:sectionGradedGOne}
  \lean{AlgebraicGeometry.Scheme.Modules.instGOneNatSectionDeg}
  \uses{def:sectionDeg, lem:moduleUnit_mathlib}
  The graded unit (a \(\mathrm{GradedMonoid.GOne}\) on
  \(m \mapsto (\mathrm{sectionDeg}\ \mathcal{L})_m\)) is the element
  \[
    1 \;\in\; (\mathrm{sectionDeg}\ \mathcal{L})_0
      = \Gamma(X, \mathcal{L}^{\otimes 0}) = \Gamma(X, \mathcal{O}_X),
  \]
  the image of \(1 \in \Gamma(X, \mathcal{O}_X)\) under the canonical
  \(\Gamma(X,\mathcal{O}_X)\)-module identification
  \(\Gamma(X, \mathcal{O}_X) = \Gamma(X, \mathcal{L}^{\otimes 0})\) coming from the
  unit module \(\mathcal{L}^{\otimes 0} = \mathbf{1}_X\)
  (\cref{lem:moduleUnit_mathlib}).
\end{definition}

\begin{lemma}[Coherence of the section multiplication]
  \label{lem:sectionMul_coherent}
  \lean{AlgebraicGeometry.Scheme.Modules.sectionsMul_one_mul,
    AlgebraicGeometry.Scheme.Modules.sectionsMul_mul_one,
    AlgebraicGeometry.Scheme.Modules.sectionsMul_mul_assoc,
    AlgebraicGeometry.Scheme.Modules.sectionsMul_mul_comm}
  \uses{def:sectionMul, def:sectionsCast, def:sectionGradedGMul,
    def:sectionGradedGOne, lem:moduleUnit_mathlib, lem:sheafTensorPow_add}
  Write \(a \cdot b\) for the degreewise section multiplication on tensor powers:
  for \(a \in \Gamma(X, \mathcal{L}^{\otimes i})\) and
  \(b \in \Gamma(X, \mathcal{L}^{\otimes j})\) it is the elementary section
  multiplication (\cref{def:sectionMul}) followed by the tensor-power comparison
  isomorphism \(\mu_{i,j}\) (\cref{lem:sheafTensorPow_add}), landing in
  \(\Gamma(X, \mathcal{L}^{\otimes (i+j)})\); and let \(1\) denote the unit of
  degree \(0\), the image of \(1 \in \Gamma(X,\mathcal{O}_X)\) under the
  identification \(\Gamma(X,\mathcal{O}_X) = \Gamma(X, \mathcal{L}^{\otimes 0})\)
  (\cref{lem:moduleUnit_mathlib}). Mirroring
  \(\mathrm{TensorPower.\{one\_mul, mul\_one, mul\_assoc\}}\) together with the
  commutativity lemma of \texttt{Mathlib.LinearAlgebra.TensorPower.Basic}, the
  multiplication satisfies the following four component-level identities. Each is
  an honest equality in a single section module, obtained by transporting one side
  along the relevant index equality via \cref{def:sectionsCast}; the subscript on
  \(\mathrm{sectionsCast}\) records that index equality.
  \begin{itemize}
    \item \emph{Left unitality.} For \(a \in \Gamma(X, \mathcal{L}^{\otimes n})\),
      \[
        \mathrm{sectionsCast}_{\,0+n=n}\,(1 \cdot a) = a .
      \]
    \item \emph{Right unitality.} For \(a \in \Gamma(X, \mathcal{L}^{\otimes n})\),
      \[
        \mathrm{sectionsCast}_{\,n+0=n}\,(a \cdot 1) = a .
      \]
    \item \emph{Associativity.} For
      \(a \in \Gamma(X, \mathcal{L}^{\otimes n_a})\),
      \(b \in \Gamma(X, \mathcal{L}^{\otimes n_b})\),
      \(c \in \Gamma(X, \mathcal{L}^{\otimes n_c})\),
      \[
        \mathrm{sectionsCast}_{\,(n_a+n_b)+n_c = n_a+(n_b+n_c)}\,
          ((a \cdot b) \cdot c) = a \cdot (b \cdot c) .
      \]
    \item \emph{Commutativity.} For
      \(a \in \Gamma(X, \mathcal{L}^{\otimes n_a})\) and
      \(b \in \Gamma(X, \mathcal{L}^{\otimes n_b})\),
      \[
        \mathrm{sectionsCast}_{\,n_a+n_b = n_b+n_a}\,(a \cdot b) = b \cdot a .
      \]
  \end{itemize}
  The transport is indispensable: the indices on the two sides of each identity
  are equal but not definitionally so, hence a priori live in distinct section
  modules, and \cref{def:sectionsCast} moves one side into the other's module to
  produce a genuine equality there.
\end{lemma}

\begin{proof}
  \uses{def:sectionMul, def:sectionsCast, lem:presheafModule_monoidal_mathlib,
    lem:presheafModule_sheafification_mathlib, lem:sheafTensorPow_add,
    cor:sheafTensorObjAssoc}
  All four identities reduce to the corresponding identities at the level of the
  tensor product presheaf, where by \cref{lem:presheafModule_monoidal_mathlib}
  evaluation at the top open is strictly monoidal, so the associator, unitors and
  symmetry act on elementary tensors by the usual formulas
  \((\sigma \otimes \tau) \otimes \upsilon \mapsto
  \sigma \otimes (\tau \otimes \upsilon)\),
  \(1 \otimes \sigma \mapsto \sigma\) and
  \(\sigma \otimes \tau \mapsto \tau \otimes \sigma\). The sheafification unit
  \(\eta\) is a natural transformation, so it commutes with the associator,
  unitors and symmetry (\cref{lem:presheafModule_sheafification_mathlib}); applying
  \(\Gamma(X,-)\) to the naturality squares transports these presheaf-level
  identities to the section multiplications, and the index transport
  \cref{def:sectionsCast} is precisely the image under \(\Gamma(X,-)\) of the
  tensor-power comparison \(\mu\) (\cref{lem:sheafTensorPow_add}) along which the
  two sides are matched.
\end{proof}

\section{Graded-ring and graded-module assembly}
\label{sec:sgr_graded_assembly}

The components and multiplications of the previous two sections are assembled
into the graded ring and graded module through Mathlib's
\emph{external-direct-sum} graded API: the relevant inputs are families of
abelian groups indexed by \(\mathbb{N}\) together with multiplication maps
landing in the index-sum component, not submodules of a pre-existing ring.

\begin{lemma}[External-direct-sum graded commutative semiring]
  \label{lem:directSum_gcommSemiring_mathlib}
  \lean{DirectSum.GCommSemiring}
  \mathlibok
  \textit{Provided by Mathlib (\texttt{Mathlib.Algebra.DirectSum.Ring}).}
  Let \(A : \mathbb{N} \to \mathrm{Type}\) be a family of additive commutative
  monoids equipped with a graded multiplication
  \(A_i \times A_j \to A_{i+j}\) and a unit \(1 \in A_0\) satisfying graded
  associativity, two-sided unitality, distributivity, and graded commutativity
  (the data of a \(\mathrm{GCommSemiring}\) on \(A\)). Then the external direct
  sum \(\bigoplus_i A_i\) is a commutative semiring, with multiplication the
  bilinear extension of the graded multiplication and identity the image of
  \(1 \in A_0\).
\end{lemma}

\begin{lemma}[External-direct-sum graded module]
  \label{lem:directSum_gmodule_mathlib}
  \lean{DirectSum.Gmodule}
  \mathlibok
  \textit{Provided by Mathlib (\texttt{Mathlib.Algebra.Module.GradedModule}).}
  Let \(A : \mathbb{N} \to \mathrm{Type}\) carry a graded semiring structure as
  in \cref{lem:directSum_gcommSemiring_mathlib} and let
  \(M : \mathbb{N} \to \mathrm{Type}\) be a family of additive commutative
  monoids with a graded scalar action \(A_i \times M_j \to M_{i+j}\) satisfying
  the graded module axioms (the data of a \(\mathrm{Gmodule}\) of \(M\) over
  \(A\)). Then the external direct sum \(\bigoplus_j M_j\) is a module over the
  semiring \(\bigoplus_i A_i\), with scalar multiplication the bilinear
  extension of the graded action.
\end{lemma}

\begin{lemma}[Graded semiring structure on the section components]
  \label{lem:sectionGradedRing_gcommSemiring}
  \lean{AlgebraicGeometry.sectionGradedRing_gcommSemiring}
  \uses{def:sheafTensorPow, lem:sheafTensorPow_add, def:sectionMul,
    lem:sectionMul_coherent, def:sectionsCast, lem:gradedMonoid_eq_of_cast,
    lem:directSum_gcommSemiring_mathlib, lem:moduleUnit_mathlib}
  % SOURCE: [Stacks Project], "Sheaves of Modules", \S "Invertible modules",
  % preceding `definition-gamma-star' (Tag 01CV) (read from
  % references/stacks-modules.tex, L4257-L4267).
  % SOURCE QUOTE:
  % "Let $(X, \mathcal{O}_X)$ be a ringed space. We can define a $\mathbf{Z}$-graded
  %  ring structure on $\bigoplus \Gamma(X, \mathcal{L}^{\otimes n})$ by mapping
  %  $s \in \Gamma(X, \mathcal{L}^{\otimes n})$ and
  %  $t \in \Gamma(X, \mathcal{L}^{\otimes m})$ to the
  %  section corresponding to $s \otimes t$ in
  %  $\Gamma(X, \mathcal{L}^{\otimes n + m})$.
  %  We omit the verification that this defines a commutative
  %  and associative ring with $1$."
  \textit{Source: [Stacks Project], "Sheaves of Modules", Tag 01CV (the graded
  ring structure on \(\bigoplus \Gamma(X,\mathcal{L}^{\otimes n})\)).}
  Let \(L_s\) be a line bundle on \(X_s\). The family of
  \(\Gamma(X_s,\mathcal{O}_{X_s})\)-modules
  \(m \mapsto \Gamma(X_s, L_s^{\otimes m})\) carries a graded commutative
  semiring structure (a \(\mathrm{GCommSemiring}\),
  \cref{lem:directSum_gcommSemiring_mathlib}) whose degree-\((m,m')\)
  multiplication is the composite
  \[
    \Gamma(X_s, L_s^{\otimes m}) \otimes_{\Gamma(X_s,\mathcal{O}_{X_s})}
      \Gamma(X_s, L_s^{\otimes m'})
    \xrightarrow{\ \cdot\ }
    \Gamma\bigl(X_s, L_s^{\otimes m} \otimes L_s^{\otimes m'}\bigr)
    \xrightarrow{\ \Gamma(\mu_{m,m'})\ }
    \Gamma\bigl(X_s, L_s^{\otimes (m+m')}\bigr),
  \]
  the section multiplication (\cref{def:sectionMul}) followed by the
  tensor-power comparison isomorphism (\cref{lem:sheafTensorPow_add}), with unit
  \(1 \in \Gamma(X_s, \mathcal{O}_{X_s}) = \Gamma(X_s, L_s^{\otimes 0})\)
  (\cref{lem:moduleUnit_mathlib}). Consequently
  \(R(X_s, L_s) = \bigoplus_{m \ge 0} \Gamma(X_s, L_s^{\otimes m})\) is a
  commutative semiring; this is the graded-ring structure underlying
  \cref{def:sectionGradedRing}.
\end{lemma}

\begin{proof}
  \uses{lem:sheafTensorPow_add, lem:sectionMul_coherent, def:sectionsCast,
    lem:gradedMonoid_eq_of_cast, def:sectionMul,
    lem:directSum_gcommSemiring_mathlib}
  The structure is assembled field-for-field as \(\mathrm{TensorPower}\) builds its
  graded semiring in \texttt{Mathlib.LinearAlgebra.TensorPower.Basic}. The
  degreewise multiplication \(\Gamma(\mu_{m,m'}) \circ (\,\cdot\,)\) of the
  statement and the degree-\(0\) unit \(1 \in \Gamma(X, \mathcal{L}^{\otimes 0})\)
  provide the graded multiplication and unit. Their graded-monoid axioms --- left
  and right unitality and associativity --- are obtained from the corresponding
  component-level identities of \cref{lem:sectionMul_coherent} by passing each
  through the bridge \cref{lem:gradedMonoid_eq_of_cast}, which converts a
  transport-mediated equality (\cref{def:sectionsCast}) into the required equality
  of dependent pairs; the graded power operation is the default iterated product
  and contributes no further data. This yields the graded-monoid layer.

  The remaining semiring axioms are the bilinearity (left and right distributivity
  and annihilation by \(0\)) of the graded multiplication and the behaviour of the
  natural-number coercion. Bilinearity is immediate: the degreewise multiplication
  \(\Gamma(\mu_{m,m'}) \circ (\,\cdot\,)\) is \(\Gamma(X,\mathcal{O}_X)\)-bilinear,
  being the composite of the bilinear section multiplication (\cref{def:sectionMul})
  with the linear comparison \(\Gamma(\mu_{m,m'})\)
  (\cref{lem:sheafTensorPow_add}), so each side vanishes on \(0\) and distributes
  over sums. The natural-number coercion sends \(n\) to \(n \cdot 1\), the
  \(n\)-fold sum of the degree-\(0\) unit. This completes the graded-semiring
  layer.

  Finally graded commutativity is the commutativity identity of
  \cref{lem:sectionMul_coherent}, again passed through
  \cref{lem:gradedMonoid_eq_of_cast}. These are precisely the data of a
  \(\mathrm{GCommSemiring}\) on \(m \mapsto \Gamma(X, \mathcal{L}^{\otimes m})\),
  so \cref{lem:directSum_gcommSemiring_mathlib} endows
  \(\bigoplus_{m \ge 0} \Gamma(X_s, L_s^{\otimes m})\) with its commutative-semiring
  structure.
\end{proof}

\begin{lemma}[Graded module structure on the twisted-section components]
  \label{lem:sectionGradedModule_gmodule}
  \lean{AlgebraicGeometry.sectionGradedModule_gmodule}
  \uses{def:sheafModuleTwist, lem:sheafTensorPow_add, def:sectionMul,
    lem:sectionMul_coherent, def:sectionsCast, lem:gradedMonoid_eq_of_cast,
    lem:directSum_gmodule_mathlib, lem:sectionGradedRing_gcommSemiring}
  % SOURCE: [Stacks Project], "Sheaves of Modules", \S "Invertible modules",
  % `definition-gamma-star' (Tag 01CV) (read from
  % references/stacks-modules.tex, L4279-L4286).
  % SOURCE QUOTE:
  % "Given a sheaf of $\mathcal{O}_X$-modules $\mathcal{F}$ we set
  %  $$
  %  \Gamma_*(X, \mathcal{L}, \mathcal{F})
  %  =
  %  \bigoplus\nolimits_{n \in \mathbf{Z}} \Gamma(X,
  %  \mathcal{F} \otimes_{\mathcal{O}_X} \mathcal{L}^{\otimes n})
  %  $$
  %  which we think of as a graded $\Gamma_*(X, \mathcal{L})$-module."
  \textit{Source: [Stacks Project], "Sheaves of Modules", Tag 01CV (the graded
  module \(\Gamma_*(X,\mathcal{L},\mathcal{F})\)).}
  Let \(\mathcal{F}_s\) be a coherent sheaf and \(L_s\) a line bundle on
  \(X_s\). The family
  \(m \mapsto \Gamma(X_s, \mathcal{F}_s \otimes L_s^{\otimes m})\)
  (\cref{def:sheafModuleTwist}) carries a graded module structure (a
  \(\mathrm{Gmodule}\), \cref{lem:directSum_gmodule_mathlib}) over the graded
  semiring of \cref{lem:sectionGradedRing_gcommSemiring}, with degree-\((m,m')\)
  action
  \[
    \Gamma(X_s, L_s^{\otimes m}) \otimes_{\Gamma(X_s,\mathcal{O}_{X_s})}
      \Gamma(X_s, \mathcal{F}_s \otimes L_s^{\otimes m'})
    \xrightarrow{\ \cdot\ }
    \Gamma\bigl(X_s, L_s^{\otimes m} \otimes
      (\mathcal{F}_s \otimes L_s^{\otimes m'})\bigr)
    \xrightarrow{\ \sim\ }
    \Gamma\bigl(X_s, \mathcal{F}_s \otimes L_s^{\otimes (m+m')}\bigr),
  \]
  the section multiplication followed by the comparison isomorphism that
  reassociates and merges the line-bundle factors (using
  \cref{lem:sheafTensorPow_add} and the symmetry of the tensor product).
  Consequently \(M(X_s, \mathcal{F}_s, L_s) = \bigoplus_{m \ge 0}
  \Gamma(X_s, \mathcal{F}_s \otimes L_s^{\otimes m})\) is a graded module over
  \(R(X_s, L_s)\); this is the graded-module structure underlying
  \cref{def:sectionGradedModule}.
\end{lemma}

\begin{proof}
  \uses{lem:sectionMul_coherent, lem:sheafTensorPow_add, def:sectionsCast,
    lem:gradedMonoid_eq_of_cast, def:sectionMul,
    lem:directSum_gmodule_mathlib, lem:sectionGradedRing_gcommSemiring}
  The action maps are the displayed composites: the section multiplication
  (\cref{def:sectionMul}) followed by the comparison isomorphism that reassociates
  and merges the line-bundle factors (\cref{lem:sheafTensorPow_add}). The
  \(\mathrm{Gmodule}\) is assembled field-for-field as the graded-action precedent
  \(\mathrm{GradedMonoid.GMulAction}\) of Mathlib, with the action grading on
  indices given by ordinary addition on \(\mathbb{N}\) (so the degree-\((m,m')\)
  action lands in degree \(m+m'\)). Its sigma-type axioms --- unitality
  \(1 \cdot x = x\) and compatibility
  \(r \cdot (r' \cdot x) = (r r') \cdot x\) with the ring multiplication --- are
  the unit and associativity coherences of the section multiplication
  (\cref{lem:sectionMul_coherent}) read with the third tensor factor carrying
  \(\mathcal{F}_s\), each converted from its transport-mediated form
  (\cref{def:sectionsCast}) into the required equality of dependent pairs by the
  bridge \cref{lem:gradedMonoid_eq_of_cast}. The bilinearity axioms of the action
  (additivity and annihilation by \(0\) in each variable) are free, the action map
  being \(\Gamma(X,\mathcal{O}_X)\)-bilinear exactly as in the proof of
  \cref{lem:sectionGradedRing_gcommSemiring}. These are the data of a
  \(\mathrm{Gmodule}\), so \cref{lem:directSum_gmodule_mathlib} produces the module
  structure of \(M(X_s, \mathcal{F}_s, L_s)\) over \(R(X_s, L_s)\).
\end{proof}
